\section*{Apendice D: Glosario de Terminos Tecnicos}

\begin{description}
    \item[Acelerometro] Sensor que mide aceleracion lineal en 3 ejes (x, y, z)

    \item[Arousale] Activacion del sistema nervioso, pasaje de reposo a alerta

    \item[Bandas EEG] Agrupaciones de frecuencias en electoencefalograma:
    Alpha (8-12 Hz), Beta (12-30 Hz), Delta (0.5-4 Hz), Theta (4-8 Hz), Gamma (30-100 Hz)

    \item[Bels] Escala logaritmica para potencia de senales (10 log10 de potencia)

    \item[BVP (Blood Volume Pulse)] Cambios en volumen sanguineo, derivado de PPG

    \item[ECG] Electrocardiograma - registro de actividad electrica del corazon

    \item[EDA (Electrodermal Activity)] Conductancia de piel, reflejo de activacion simpatica

    \item[EEG] Electroencefalograma - registro de actividad electrica cerebral

    \item[Fotopletismograma (PPG)] Sensor optico que mide cambios en volumen sanguineo

    \item[Giroscopo] Sensor que mide velocidad angular/rotacion

    \item[HR (Heart Rate)] Numero de latidos por minuto

    \item[HRV (Heart Rate Variability)] Variacion en intervalos entre latidos consecutivos (SDNN en ms)

    \item[SDNN] Standard Deviation of Normal-to-Normal intervals - medida de HRV

    \item[Simpatico] Rama del sistema nervioso autonomo responsable de "lucha o huida"

    \item[SNA] Sistema Nervioso Autonomo - controla funciones involuntarias

    \item[Parasimpatico] Rama del SNA responsable de "descanso y digestion"

    \item[Wearable] Dispositivo portatil/llevable incorporado en ropa o accesorios
\end{description}
